% ============================
%   4. Architecture
% ============================
\section{System architecture}
\label{sec:architecture}

\subsection{High-level overview}

\nidsname{} is a two-tier application: a Python backend that performs
capture, detection and persistence; and a single-page React frontend
that consumes its REST and WebSocket APIs. The two communicate over
HTTP/JSON for queries and over two dedicated WebSocket channels
(\code{/ws/alerts} and \code{/ws/traffic}) for real-time push.
Figure~\ref{fig:dataflow} sketches the packet life-cycle.

\begin{figure}[h]
\centering
\fbox{
\parbox{0.92\textwidth}{\centering\ttfamily
\small
NIC / Simulator
\;$\rightarrow$\; Sniffer
\;$\rightarrow$\; Dispatcher
\;$\rightarrow$\; [ Signature engine \;+\; ML detector ]\\[0.3em]
\;$\rightarrow$\; merge \;$\rightarrow$\; DB (Alert, PcapRecord, TrafficEvent)\\[0.3em]
\;$\rightarrow$\; per-alert JSONL ring on disk \;+\; WebSocket broadcast \;$\rightarrow$\; Browser
}}
\caption{End-to-end packet flow from capture to UI.}
\label{fig:dataflow}
\end{figure}

\subsection{Component breakdown}

\subsubsection{Capture}

Two implementations share a common asynchronous callback signature:
\begin{itemize}
    \item \code{LiveSniffer} uses Scapy~\cite{scapy} on a real interface.
        It runs on a dedicated thread and posts each parsed
        \code{Packet} to the main event loop via
        \code{asyncio.run\_coroutine\_threadsafe}.
    \item \code{TrafficSimulator} emits a steady benign baseline of
        24--40 packets per second and injects an attack burst every
        five seconds (Nmap-like scan, SSH brute-force, ICMP flood, SQL
        injection, command injection, reverse shell). It exists so
        that the application can be demonstrated on a laptop without
        the elevated privileges required by raw packet capture, and so
        that the ML detector always has labelled-ish data to react to.
\end{itemize}

\subsubsection{Dispatcher}

The \code{Dispatcher} is the central coordination point. For every
received \code{Packet} it:
\begin{enumerate}
    \item Increments the packet counter and the sliding-window packets-per-second
        meter.
    \item Pushes the packet onto a fixed-size ring buffer (\code{deque}
        of size 200), which acts as the source for forensic snapshots.
    \item Publishes a compact event onto the \code{traffic\_bus} so that
        any subscribed WebSocket clients (i.e. the dashboard) see live
        throughput.
    \item Optionally persists a \code{TrafficEvent} record (sampled at
        $1/5$ packets to keep the events table small during demos).
    \item Calls the signature engine and the ML detector and merges
        their verdicts into zero-or-more alert candidates.
    \item For each candidate, writes a JSONL ring-tail to disk, creates
        \code{Alert} and \code{PcapRecord} rows, and broadcasts the
        alert on \code{alert\_bus}.
\end{enumerate}

\subsubsection{Signature engine}

The engine compiles each enabled DB row into an in-memory
\code{CompiledRule}: a struct of pre-compiled regular expressions,
port ranges, flag masks, and (for threshold rules) a sliding-window
counter keyed by source IP. Compilation happens once at startup and
again whenever the rules table is mutated through the API (hot reload).
For every incoming packet, the engine iterates all enabled rules and
yields one alert-candidate dictionary per match.

\subsubsection{ML detector}

The \code{MLDetector} class wraps both algorithms behind a common
\code{score(pkt) $\rightarrow$ (label, score)} interface. State (current
model, feature names, per-source last-seen timestamp) is protected by a
mutex so that training can run concurrently with inference.
Persistence uses \code{joblib}; activating a model loads the saved
pickle and flips the singleton's active flag.

\subsubsection{Forensics layer}

When an alert is created, \code{write\_packet\_log()} serialises the
last 50 entries of the dispatcher's ring buffer as JSON Lines into a
file named after the alert ID. The path is stored on the \code{Alert}
row, and a separate \code{PcapRecord} row mirrors it for table-driven
browsing in the UI. Download is served by an authenticated FastAPI
endpoint that performs a path check and streams the file with
\code{FileResponse}.

\subsubsection{REST + WebSocket API}

The HTTP surface is organised by tag in FastAPI:
\code{/api/auth}, \code{/api/alerts}, \code{/api/rules},
\code{/api/dashboard}, \code{/api/forensics}, \code{/api/ml},
\code{/api/health}. All resource routes are protected by
\code{Depends(get\_current\_user)}, which decodes a JWT and reloads the
caller from the database. Two WebSocket endpoints push real-time data:
\code{/ws/alerts} (each new alert as a single message) and
\code{/ws/traffic} (batched traffic samples at $\approx 4$\,Hz to keep
the dashboard chart smooth without overwhelming the network).

\subsubsection{Frontend}

The dashboard is a Vite~\cite{vite}-bundled React application styled
with TailwindCSS. State is held mostly in React Context (auth,
live alerts, live traffic); pages such as \code{Alerts} and \code{Rules}
fetch via authenticated REST and reconcile against the live WebSocket
stream client-side. The visual identity (deep navy with cyan accents) is
applied through a Tailwind theme to keep the look consistent across all
pages and the new modal components introduced in Section~\ref{sec:current_state}.

\subsection{Data model}

The application uses five SQLAlchemy ORM tables:
\begin{description}
    \item[User] account credentials and role (admin/analyst).
        Passwords are hashed with bcrypt~\cite{bcrypt}.
    \item[Rule] one row per signature; the \code{match} field is a JSON
        blob whose shape mirrors Listing~\ref{lst:rule}.
    \item[Alert] the central record. Carries severity, category,
        detection method (signature/ml/hybrid), ML score, source/destination
        5-tuple, status (new/triaged/closed), and the path to its
        associated PCAP file.
    \item[PcapRecord] denormalised view of the forensic file for the UI
        Forensics table; references \code{Alert} via foreign key.
    \item[TrafficEvent] sampled record of traffic for historical
        analysis.
    \item[MLModel] persisted metadata for trained models (algorithm,
        path, metrics, feature importance, active flag).
\end{description}

The schema is created with \code{Base.metadata.create\_all()} on first
boot; an Alembic migration suite is planned for the final thesis to
support schema evolution without losing data.

\subsection{One-click launcher}

The launcher (\code{run.py}) is a deliberate design choice rather than
an afterthought. It performs, in order: Python version check (refuses
$\geq 3.14$ because scientific packages lack Windows wheels for very
recent CPython releases, which would force a Visual~C++ build);
creation of \code{backend/.venv}; dependency installation against a
pinned \code{requirements.txt}; frontend \code{npm install};
parallel launch of \code{uvicorn} and \code{vite dev}; browser
auto-open; supervised waiting on Ctrl+C with cooperative shutdown of
both children. The wrappers \code{run.bat} (Windows) and \code{run.sh}
(macOS/Linux) merely pick the best available Python interpreter and
delegate to \code{run.py}.
